\section{維持管理を点検から次の判断へつなぐ}\label{theory:safety:maintenance}
社会資本の機能を持続させるには、状態を点検し、診断し、必要な措置を選び、
その根拠と結果を記録して次の周期へ反映する。
予防的・計画的な修繕は、劣化が大きくなってから一斉更新する負担を避け、
長期の費用の縮減・平準化と利用者に必要な機能の維持へつなげる。
計画を作れば必ず費用が一定になる、という保証ではない。\cite{eng52_infra2014}

平成26年の行動計画は施設分野別の基準・取組を含む。
施設の重要度、状態、材料、環境と分野ごとの最低基準を踏まえて実施するのであり、
全インフラの頻度と内容を一つへ固定するという意味ではない。
分野内での統一した健全度の尺度と、全施設で同じ点検をするという全称を分ける。
構造諸元、点検履歴、修繕記録、損傷位置・変化を蓄積し、
データが不足している状態を把握することも診断の仕事である。

\subsection{新技術と人材・責任を結び付ける}
ロボット、非破壊検査、ICTで効率・観測範囲を改善できるが、
測定の限界・見落とし・適用可能性を評価する能力が必要になる。
施設分野・業務ごとの必要能力を明確にし、担い手を確保・育成する。
資格の保有だけで全分野の経験と判断を代替しない。\cite{eng52_infra2014}

\subsection{行政責任と市民参加は両立する}
管理者の安全確保の責任を、地域住民・民間の協力を一律禁止する根拠にはしない。
国土交通省の平成12年白書にも、行政との合意に基づく道路・河川の清掃等の協働が記録されている。
誰が何を行い、どこから専門点検・管理者の判断が必要になるかを明確にする。
清掃への参加の事実から危険作業や構造診断を無条件に住民へ委任できるとは推論しない。\cite{eng52_adopt}
